\documentclass[11pt,a4paper]{article}

\usepackage[spanish]{babel}
\usepackage[utf8]{inputenc}
\usepackage[T1]{fontenc}
\usepackage{lmodern}
\usepackage{microtype}
\usepackage[margin=2.5cm]{geometry}
\usepackage{amsmath}
\usepackage{amssymb}
\usepackage{listings}
\usepackage{xcolor}
\usepackage{hyperref}
\usepackage{enumitem}
\usepackage{titlesec}
\usepackage{fancyhdr}

\hypersetup{
  colorlinks=true,
  linkcolor=blue!50!black,
  urlcolor=blue!50!black,
  citecolor=blue!50!black,
}

\definecolor{codebg}{RGB}{245,245,247}
\definecolor{codekey}{RGB}{0,90,156}
\definecolor{codestr}{RGB}{160,80,30}
\definecolor{codecom}{RGB}{120,120,120}

\lstset{
  basicstyle=\ttfamily\footnotesize,
  backgroundcolor=\color{codebg},
  keywordstyle=\color{codekey}\bfseries,
  stringstyle=\color{codestr},
  commentstyle=\color{codecom}\itshape,
  showstringspaces=false,
  breaklines=true,
  frame=single,
  rulecolor=\color{black!20},
  numbers=left,
  numberstyle=\tiny\color{black!50},
  numbersep=8pt,
  xleftmargin=12pt,
  framexleftmargin=4pt,
  captionpos=b,
}

\titleformat{\section}{\Large\bfseries}{\thesection.}{0.6em}{}
\titleformat{\subsection}{\large\bfseries}{\thesubsection}{0.6em}{}

\pagestyle{fancy}
\fancyhf{}
\rhead{\small PortfolioLab · Guía técnica}
\lhead{\small TFG · Rubén Moreno}
\cfoot{\thepage}

\title{\textbf{PortfolioLab}\\\large Guía técnica del código\\\small (documento interno de trabajo, base de la memoria del TFG)}
\author{Rubén Moreno Fernández}
\date{Mayo 2026}

\begin{document}
\maketitle
\thispagestyle{empty}

\begin{abstract}
\noindent Este documento describe módulo a módulo qué hace cada parte del código de la aplicación PortfolioLab, las decisiones de diseño que hay detrás, las fórmulas matemáticas implementadas y las referencias bibliográficas relevantes. Está pensado como guía de estudio personal antes de escribir la memoria definitiva del TFG y como referencia rápida durante la defensa.
\end{abstract}

\tableofcontents
\newpage

% =============================================================================
\section{Visión general y arquitectura}
% =============================================================================

PortfolioLab es una aplicación web full-stack que combina tres ideas que en un TFG raramente aparecen juntas:

\begin{enumerate}
  \item Un \textbf{menú de cinco métodos de optimización} de carteras (Markowitz, mínima varianza, Risk Parity, mínimo CVaR, Equal Weight) que comparten el mismo marco de matrices de covarianzas pero responden a filosofías de gestión distintas.
  \item Un módulo de \textbf{planificación financiera personal} estilo \textit{robo-advisor} (cuánto invertir, cómo distribuirlo) que conecta la teoría con la situación real del inversor.
  \item Un análisis del impacto del \textbf{IRPF español} (base imponible del ahorro) y del diferimiento fiscal de los Fondos de Inversión, raramente tratado en TFGs cuantitativos.
\end{enumerate}

La aplicación se compone de tres capas:

\begin{description}[leftmargin=1.4cm,labelwidth=1.2cm]
  \item[Frontend] React 19 + TypeScript + Vite. Páginas en \texttt{src/pages/}, componentes reutilizables en \texttt{src/components/}, lógica pura en \texttt{src/services/}.
  \item[Backend] FastAPI + SQLAlchemy. Endpoints en \texttt{app/api/}, lógica financiera en \texttt{optimizer/} y \texttt{backtesting/}, ETL en \texttt{etl/}.
  \item[Base de datos] PostgreSQL con migraciones gestionadas por Alembic.
\end{description}

\noindent El flujo típico de un usuario es:

\begin{enumerate}
  \item Se autentica (JWT).
  \item En el \textit{Planificador} introduce su situación financiera y obtiene una recomendación de capital invertible y \textit{asset allocation}.
  \item Aplica esa recomendación al \textit{Optimizador}, elige un método y obtiene una cartera óptima.
  \item Guarda la cartera y la analiza en el \textit{X-Ray} (sectores, países, dividendos, clase de activo).
  \item Ejecuta un \textit{backtest} histórico con rebalanceo y comisiones, o una \textit{proyección Monte Carlo} a futuro.
  \item Compara varias carteras en el \textit{Comparador} o estudia el impacto fiscal en \textit{Fiscalidad}.
\end{enumerate}

% =============================================================================
\section{Capa de datos: ETL y caché}
% =============================================================================

\subsection{\texttt{etl/cache.py} — caché TTL en memoria}

Yfinance impone un \textit{rate-limit} implícito y descargar 2--20 años de cotizaciones para 8 tickers tarda entre 8 y 15 segundos. Llamar al optimizador o al backtest varias veces seguidas saturaba la demo, así que implementé una caché TTL en memoria sin dependencias externas.

\paragraph{Diseño}
Una clase \texttt{TTLCache} con un \texttt{OrderedDict} interno (LRU implícito al mover entradas al final) y un \texttt{threading.RLock} para ser \textit{thread-safe} en el escenario en que FastAPI sirve varias peticiones concurrentes. El decorador \texttt{ttl\_cache} envuelve funciones puras y construye la clave a partir de \texttt{(nombre\_funcion, args, kwargs)}.

\paragraph{Singletons}
Tres instancias compartidas en toda la aplicación:

\begin{itemize}
  \item \texttt{prices\_cache} con TTL 6 h y máximo 256 entradas (precios diarios, cambian poco una vez cierra la sesión).
  \item \texttt{info\_cache} con TTL 1 h y máximo 512 (metadatos del activo, sector, market cap).
  \item \texttt{search\_cache} con TTL 15 min y máximo 256 (búsqueda de tickers en Yahoo Finance).
\end{itemize}

\paragraph{Métricas observables}
Cada caché lleva un contador de \textit{hits}, \textit{misses} y calcula el \textit{hit rate}. El endpoint \texttt{GET /assets/cache-stats} expone estas métricas, lo que permite demostrar en defensa que el sistema está realmente acelerando peticiones (en pruebas, las llamadas repetidas pasan de 6\,000 ms a 0 ms).

\subsection{\texttt{etl/market\_data.py} — conector de datos de mercado}

\paragraph{Estrategia de fuentes}
Se intenta primero \texttt{yfinance} (rápido, gratis, datos completos). Si falla, se cae a \texttt{Alpha Vantage} como respaldo (requiere API key gratuita). Si ambos fallan, se lanza \texttt{DataSourceError} con un mensaje claro.

\paragraph{Decisión clave: \texttt{auto\_adjust=True}}
Por defecto se descargan precios \textbf{ajustados por dividendos y splits} (\textit{total return}). Esto significa que cuando el optimizador o el backtest calculan retornos, ya están incluyendo los dividendos reinvertidos. Es el estándar académico y comercial, aunque conviene exponerlo explícitamente al usuario.

\paragraph{Método \texttt{get\_historical\_full()}}
Para el módulo de dividendos del backtest necesitamos \textbf{descomponer} el retorno en precio puro vs dividendos. Este método descarga con \texttt{auto\_adjust=False}, lo que devuelve dos columnas:
\begin{itemize}
  \item \texttt{Close}: precio ajustado por splits pero \textbf{no} por dividendos.
  \item \texttt{Adj Close}: precio ajustado por splits \textbf{y} dividendos (total return).
\end{itemize}
Y adicionalmente la columna \texttt{Dividends} con el dividendo bruto pagado cada día. Con esto el motor de backtest puede separar la parte de retorno que viene del precio (con \texttt{Close}) de la parte que viene de los dividendos (la diferencia entre ambas series).

% =============================================================================
\section{Optimizador de carteras}
% =============================================================================

El núcleo cuantitativo de la aplicación está en \texttt{backend/optimizer/markowitz.py}. La clase \texttt{MarkowitzOptimizer} expone cinco métodos de optimización y dos análisis adicionales (frontera eficiente y frontera de Pareto). Todos comparten la misma matriz de covarianzas anualizada.

\subsection{Inicialización}

Al construir la clase con un periodo (típicamente 2 años de histórico) se descargan los precios de todos los tickers, se intersecta el calendario y se calcula:

\begin{equation}
\Sigma = \mathrm{cov}(R_t)\cdot 252,\qquad \mu = \mathbb{E}[R_t]\cdot 252,
\end{equation}

donde $R_t$ son los retornos diarios y 252 es el factor de anualización (días de cotización al año).

\paragraph{Validación añadida tras revisión}
Si la intersección temporal queda vacía (un activo demasiado nuevo respecto a los demás) o si algún ticker se pierde silenciosamente por falta de datos, se lanza un error explícito en lugar de continuar con la cartera incompleta.

\subsection{Método 1 — Markowitz: maximización del Sharpe Ratio}

Es el método clásico (Markowitz, 1952). Se busca el vector de pesos $w$ que maximiza:

\begin{equation}
\mathrm{Sharpe}(w) = \frac{w^\top\mu - r_f}{\sqrt{w^\top \Sigma w}}
\end{equation}

con $r_f = 2\,\%$ (tasa libre de riesgo), sujeto a:
\begin{align}
\sum w_i &= 1 \\
0 \le w_i &\le 1 \quad \text{(no se permiten posiciones cortas)} \\
\sqrt{w^\top \Sigma w} &\le \sigma_{\max} \quad \text{(volatilidad máxima configurada)}
\end{align}

\paragraph{Estrategia de optimización}
Se usa SLSQP (Sequential Least Squares Programming) de \texttt{scipy.optimize.minimize} con \textbf{multi-start}: el solver se lanza desde 21 puntos iniciales distintos (uno equiponderado más 20 muestreados de una distribución de Dirichlet con semilla fija para reproducibilidad). Esto evita quedarse atrapado en mínimos locales.

\paragraph{\textit{Fallback} robusto}
Si \textit{ningún} arranque converge (puede pasar con $\sigma_{\max}$ muy restrictivo), se intenta calcular la cartera de mínima varianza global. Si su volatilidad resultante supera la restricción del usuario, se lanza un error con un mensaje claro indicando la volatilidad mínima alcanzable.

\subsection{Método 2 — Mínima varianza global}

Mismo problema pero con función objetivo:

\begin{equation}
\min_w\ w^\top \Sigma w
\end{equation}

sin restricción de retorno objetivo. Devuelve la cartera más estable posible. Académicamente interesante porque empíricamente suele tener mejor \textit{out-of-sample} performance que la cartera de máximo Sharpe (los retornos esperados $\mu$ son notoriamente inestables, mientras que $\Sigma$ es algo más estable en el tiempo).

\subsection{Método 3 — Risk Parity (Equal Risk Contribution)}

Implementación del método propuesto por Maillard, Roncalli y Teïletche (2010). La idea es repartir \textit{el riesgo} entre los activos en lugar de \textit{el peso}. La contribución al riesgo de cada activo es:

\begin{equation}
\mathrm{RC}_i = w_i \cdot \frac{(\Sigma w)_i}{\sqrt{w^\top \Sigma w}}
\end{equation}

y queremos $\mathrm{RC}_i = \mathrm{RC}_j$ para todo par $(i,j)$. En la práctica se minimiza la varianza de las contribuciones marginales no normalizadas $w_i \cdot (\Sigma w)_i$:

\begin{equation}
\min_w \sum_{i,j} \big[ w_i (\Sigma w)_i - w_j (\Sigma w)_j \big]^2
\end{equation}

Es la filosofía detrás de fondos como el famoso \textit{All Weather} de Bridgewater (Ray Dalio) y se considera el estado del arte en gestión institucional moderna.

\subsection{Método 4 — Mínimo CVaR}

\textit{Conditional Value at Risk} o \textit{Expected Shortfall} al nivel $\alpha=95\,\%$. Mientras que el VaR responde a "¿cuánto puedo perder con probabilidad 95\,\%?", el CVaR responde a "y cuando la cosa se pone mal, ¿cuánto pierdo en promedio?":

\begin{equation}
\mathrm{CVaR}_\alpha(w) = \mathbb{E}\big[\,L \;|\; L \ge \mathrm{VaR}_\alpha\,\big]
\end{equation}

donde $L = -w^\top R$ son las pérdidas. La formulación clásica de optimización es la de Rockafellar y Uryasev (2000). En este TFG se implementa una versión empírica directa: para cada vector $w$ candidato se calculan los retornos diarios históricos $\{w^\top R_t\}$, se identifican los del peor 5\,\% y se promedia la pérdida sobre ellos. Después se minimiza esa cantidad mediante SLSQP con multi-start.

\paragraph{Por qué importa}
La varianza penaliza por igual subidas y bajadas. El CVaR penaliza solo el lado malo. Es una medida \textbf{coherente} de riesgo (Artzner et al., 1999) y es la métrica regulatoria recomendada por Basilea III en sustitución del VaR.

\subsection{Método 5 — Equal Weight (1/N)}

Sin optimización: $w_i = 1/N$ para todos los activos. Es un \textit{baseline} ingenuo. DeMiguel, Garlappi y Uppal (2009) demostraron empíricamente que esta cartera tan tonta no es estadísticamente peor que muchos métodos sofisticados \textit{out of sample}, debido al ruido en los parámetros estimados. Lo incluyo como referencia honesta.

\subsection{Métricas de calidad de la cartera}

Tras resolver cualquiera de los cinco métodos se devuelven seis métricas que permiten comparar los resultados con criterio:

\begin{itemize}
  \item \textbf{Retorno esperado} $w^\top\mu$ y \textbf{volatilidad} $\sqrt{w^\top \Sigma w}$.
  \item \textbf{Sharpe Ratio}.
  \item \textbf{Diversification Ratio} (Choueifaty y Coignard, 2008): $\mathrm{DR}(w) = \frac{w^\top \sigma}{\sqrt{w^\top \Sigma w}}$, donde $\sigma$ es el vector de volatilidades individuales. Vale 1 si los activos están perfectamente correlacionados, mayor cuanto más diversificación efectiva existe.
  \item \textbf{Concentración HHI} (Herfindahl-Hirschman): $\mathrm{HHI} = \sum w_i^2$, entre $1/N$ (uniforme) y $1$ (concentrada en un solo activo).
  \item \textbf{Activos efectivos}: $1/\mathrm{HHI}$, equivalente a "cuántos activos con peso uniforme darían la misma concentración".
  \item \textbf{CVaR 95\,\% diario}: pérdida media en el peor 5\,\% de días (ya descrito).
\end{itemize}

\subsection{Frontera Eficiente de Markowitz}

Se calculan 50 puntos: para cada nivel de retorno objetivo entre $\min(\mu)$ y $\max(\mu)$ se resuelve la minimización de varianza con esa restricción de retorno. El resultado es la curva clásica retorno vs volatilidad. La cartera de máximo Sharpe se sitúa en el punto donde la recta tangente desde $r_f$ toca la frontera (la \textit{Capital Market Line}).

\subsection{Frontera de Pareto Sharpe-vs-Sortino}

Innovación introducida por sugerencia del tutor. Se exploran 20 puntos a lo largo del parámetro $\theta \in [0,1]$:

\begin{equation}
\max_w\ \theta \cdot \mathrm{Sharpe}(w) + (1-\theta) \cdot \mathrm{Sortino}(w)
\end{equation}

donde el ratio de Sortino es como el Sharpe pero usando solo la \textit{volatilidad bajista} (desviación típica de los retornos negativos) en el denominador. Tras obtener los candidatos se filtran los \textit{dominados} (un punto $P$ está dominado si existe $Q$ con $\mathrm{Sharpe}(Q) \ge \mathrm{Sharpe}(P)$ y $\mathrm{Sortino}(Q) \ge \mathrm{Sortino}(P)$ y al menos una desigualdad estricta). Lo que queda es el frente de Pareto en el espacio de las dos métricas.

\paragraph{Lectura}
$\theta=1$ recupera la cartera de máximo Sharpe pura. $\theta=0$ optimiza únicamente el Sortino (al inversor solo le importa el riesgo bajista). Valores intermedios son carteras de compromiso. Visualmente se muestra como un \textit{scatter plot} con gradiente de color según $\theta$.

% =============================================================================
\section{Backtesting histórico}
% =============================================================================

El motor de \texttt{backend/backtesting/engine.py} simula el comportamiento histórico de una cartera con pesos fijos durante un periodo determinado, calculando métricas de riesgo y comparándola con un benchmark (SPY).

\subsection{Construcción de la serie de la cartera}

Hay dos modos:

\paragraph{Buy and Hold (sin rebalanceo)}
Se ponderan los retornos diarios de cada activo por su peso objetivo y se compone:

\begin{equation}
V_t = V_{t-1}\cdot\bigg(1 + \sum_i w_i^* r_{i,t}\bigg)
\end{equation}

con $V_0 = 100$. Es la versión más rápida pero ignora el \textit{drift} de los pesos: si una acción sube mucho su peso efectivo crece, y eso desvía la cartera del objetivo del usuario.

\paragraph{Con rebalanceo periódico}
Se simula día a día. Los pesos se dejan derivar con los retornos diarios y, en cada fecha de rebalanceo (mensual, trimestral, semestral o anual), se recalculan al objetivo aplicando un coste:

\begin{equation}
\mathrm{turnover} = \tfrac{1}{2}\sum_i |w_i^{\text{actual}} - w_i^*|,\qquad \mathrm{coste} = \mathrm{turnover}\cdot 2\cdot c
\end{equation}

donde $c$ es la comisión por operación (en porcentaje). El factor 2 refleja que rebalancear implica vender y comprar.

\subsection{Métricas calculadas}

Sobre la serie temporal se calculan:

\begin{itemize}
  \item \textbf{Rentabilidad acumulada}: $V_T/V_0 - 1$.
  \item \textbf{Retorno anualizado}: $(V_T/V_0)^{252/T} - 1$.
  \item \textbf{Volatilidad anualizada}: $\sigma\sqrt{252}$.
  \item \textbf{Sharpe Ratio}, \textbf{Sortino Ratio} (denominador es la volatilidad bajista anualizada), \textbf{Calmar Ratio} (retorno anualizado entre $|\mathrm{MaxDD}|$).
  \item \textbf{Maximum Drawdown}: la peor caída desde un máximo histórico hasta el siguiente mínimo.
  \item \textbf{Beta} respecto a SPY: $\mathrm{cov}(R_p, R_b)/\mathrm{var}(R_b)$.
\end{itemize}

\subsection{Análisis de crisis históricas}

Se aplican las mismas métricas, restringidas a tres periodos de estrés:

\begin{itemize}
  \item \textbf{COVID}: 19/02/2020 -- 23/03/2020.
  \item \textbf{Lehman}: 01/09/2008 -- 09/03/2009.
  \item \textbf{Corrección 2022}: 03/01/2022 -- 12/10/2022.
\end{itemize}

Sirve para que el inversor pueda contestar "¿cómo se habría comportado mi cartera durante el pánico financiero más cercano?".

\subsection{Descomposición precio vs dividendos}

Cuando los datos vienen vía \texttt{get\_historical\_full()}, se ejecutan dos backtests en paralelo: uno con la serie \textit{total return} (incluye dividendos reinvertidos) y otro con la serie \textit{precio puro}. La diferencia entre ambos retornos acumulados es la \textit{contribución de los dividendos}. Adicionalmente, agrupando los dividendos brutos por año, se calcula el flujo de ingresos anuales en € por cada 100€ invertidos.

% =============================================================================
\section{Simulación Monte Carlo}
% =============================================================================

El módulo \texttt{backend/optimizer/montecarlo.py} proyecta el comportamiento futuro de una cartera mediante simulación Monte Carlo.

\subsection{Decisión metodológica: bootstrap histórico}

Existen dos familias clásicas de simulación de retornos:

\begin{description}[leftmargin=1.4cm]
  \item[Modelo paramétrico Gaussiano] Asume $R_t \sim \mathcal{N}(\mu, \Sigma)$ y muestrea de esa distribución. Es el método más común en TFGs por su simplicidad. \textbf{Problema}: los retornos reales tienen colas más gordas de lo normal, asimetría negativa y \textit{cluster} de volatilidad. Asumir normalidad subestima la probabilidad de eventos extremos.
  \item[Bootstrap histórico] Sortea con reemplazo retornos diarios reales del histórico. Conserva las propiedades empíricas de la distribución (colas, asimetría, autocorrelaciones parciales).
\end{description}

\noindent En este TFG se eligió bootstrap histórico. La principal limitación que reconozco explícitamente al usuario es que no genera escenarios fuera de la muestra: el peor día simulado posible es el peor día observado.

\subsection{Algoritmo}

\begin{enumerate}
  \item Calcular el vector de retornos diarios de la cartera: $r_t = w^\top R_t$ a partir del histórico.
  \item Generar una matriz $S\times D$ de índices aleatorios (con reemplazo) en el rango $[0, T)$, donde $S$ es el número de simulaciones y $D = 252 \cdot a\tilde{n}os$.
  \item Aplicar la fórmula del interés compuesto vectorialmente para obtener $S$ trayectorias de capital.
  \item Calcular percentiles 5, 25, 50, 75 y 95 sobre cada paso temporal.
\end{enumerate}

\paragraph{Rendimiento}
La operación se realiza íntegramente en NumPy con \texttt{np.random.default\_rng(42)} para reproducibilidad. 5\,000 simulaciones de 10 años se ejecutan en aproximadamente 270\,ms.

\subsection{Resultados devueltos}

Para visualización se hace muestreo a paso mensual (en lugar de diario, lo cual saturaría el frontend). Adicionalmente se calculan:

\begin{itemize}
  \item Estadísticos del valor final: media, mediana, desviación, percentiles.
  \item CAGR percentilado: $\mathrm{CAGR} = (V_T/V_0)^{1/T} - 1$ por simulación, con percentiles.
  \item Probabilidades clave: $P(\text{pérdida})$, $P(\text{doblar capital})$, $P(\text{triplicar capital})$.
  \item Histograma logarítmico del valor final (la cola derecha es muy larga, escala lineal pierde resolución).
\end{itemize}

% =============================================================================
\section{Análisis fiscal español}
% =============================================================================

El módulo \texttt{backend/optimizer/fiscalidad.py} cuantifica el impacto del IRPF de la base del ahorro y la ventaja del diferimiento fiscal característico del marco español.

\subsection{Tramos IRPF de la base del ahorro (2025)}

\begin{center}
\begin{tabular}{l|c}
Base imponible del ahorro & Tipo \\\hline
0 -- 6\,000\,€         & 19\,\% \\
6\,000 -- 50\,000\,€   & 21\,\% \\
50\,000 -- 200\,000\,€ & 23\,\% \\
200\,000 -- 300\,000\,€ & 27\,\% \\
$>$ 300\,000\,€         & 28\,\% \\
\end{tabular}
\end{center}

La función \texttt{calcular\_irpf\_ahorro(base)} aplica los tramos progresivamente y devuelve el importe total a pagar más el desglose por tramo y el tipo efectivo.

\subsection{Vehículo 1: ETF distributivo}

Modelo: el inversor cobra dividendos cada año, se le aplica una retención del 19\,\% en el momento del pago (a cuenta del IRPF), y reinvierte el neto. Al rescatar, paga IRPF sobre la plusvalía.

\paragraph{Detalle clave del coste fiscal de adquisición}
Los netos reinvertidos cada año amplían la \textit{base de coste} del inversor a efectos fiscales. Si tiene la cartera 20 años y reinvierte X € en netos, su coste fiscal es \textit{capital inicial + X}, no solo capital inicial. Por tanto la plusvalía gravable al rescatar es:

\begin{equation}
\mathrm{plusvalia} = V_T - \big(C_0 + \textstyle\sum_t \text{neto}_t\big)
\end{equation}

\paragraph{Aviso}
Una primera implementación del módulo no incluía esa corrección, lo que provocaba doble tributación: el dividendo se gravaba como retención y, al ampliarse el valor de la cartera, también como plusvalía al rescatar. Tras revisión se corrigió. La cifra resultante de "ahorro por diferimiento" pasó del 6.5\,\% (errónea) al 1\,\% aproximado, mucho más realista.

\subsection{Vehículo 2: Fondo de Inversión / ETF acumulativo}

Los dividendos se reinvierten dentro del producto sin pasar por el inversor: cero retenciones intermedias, todo el rendimiento sigue componiéndose. El IRPF solo se paga al rescatar, sobre la plusvalía total. Es el famoso \textbf{diferimiento fiscal} del marco regulatorio español (Ley 35/2006), peculiar y muy ventajoso a largo plazo.

\subsection{Comparativa}

La función \texttt{comparar()} ejecuta ambos escenarios con los mismos parámetros (capital, años, CAGR esperado, yield de dividendos) y devuelve el ahorro absoluto y porcentual del Fondo respecto al ETF distributivo. Para horizontes largos y yields altos, el diferimiento puede ser determinante.

\paragraph{Disclaimer académico}
La simulación es \textbf{simplificada}: no contempla compensación de pérdidas patrimoniales, regla de los 2 meses, casos transfronterizos, fiscalidad foral (País Vasco/Navarra) ni doctrina más reciente. La retención del 19\,\% se trata como tipo final cuando la base anual no supera 6\,000\,€. No constituye asesoramiento fiscal.

% =============================================================================
\section{Planificador financiero (\textit{robo-advisor})}
% =============================================================================

La página \texttt{Planificador} y su lógica pura en \texttt{frontend/src/services/planner.ts} implementan un asesor financiero personal simplificado.

\subsection{Inputs del usuario}

\begin{itemize}
  \item Liquidez disponible (€).
  \item Gastos mensuales (€).
  \item Meses deseados de fondo de emergencia (3, 6 o 12).
  \item Edad.
  \item Horizonte de inversión (corto $<$2y, medio 2--5y, largo 5--10y, muy largo $>$10y).
  \item Perfil de riesgo (conservador, moderado, agresivo).
\end{itemize}

\subsection{Cálculo de la asignación recomendada}

\paragraph{Fondo de emergencia}
$\mathrm{emergencia} = \mathrm{gastos} \times \mathrm{meses}$. Se reserva en cuenta corriente o cuenta remunerada y queda fuera del capital invertible.

\paragraph{Capital invertible}
$\mathrm{invertible} = \max(0, \mathrm{liquidez} - \mathrm{emergencia})$.

\paragraph{Porcentaje de equity (regla del 110)}
\begin{equation}
\%\mathrm{equity} = 110 - \mathrm{edad} + \Delta_{\text{perfil}} + \Delta_{\text{horizonte}}
\end{equation}

con $\Delta_{\text{perfil}} \in \{-10, 0, +10\}$ pp según conservador/moderado/agresivo y $\Delta_{\text{horizonte}}$ que penaliza el corto plazo y bonifica el largo plazo. Adicionalmente se aplica un \textit{cap} por horizonte (a corto plazo, el equity no puede superar el 30\,\% por mucho que la fórmula lo sugiera).

\paragraph{Reparto de la parte no equity}
El resto se distribuye entre bonos, real estate y commodities. La parte de \textit{alternativos} (real estate + commodities) se ajusta al perfil: el conservador casi todo en bonos, el agresivo más diversificado entre alternativos.

\subsection{Notas contextuales automáticas}

Seis reglas de aviso disparan mensajes según la situación: si el fondo de emergencia consume más del 70\,\% de la liquidez, si el horizonte corto choca con un porcentaje de equity alto, si la edad y la regla del 110 entran en contradicción, etc.

\subsection{Integración con el optimizador}

Al pulsar "Aplicar al Optimizador", se persiste el capital invertible y la tolerancia al riesgo en el perfil del usuario, y se almacena la \textit{asset allocation} recomendada en \texttt{localStorage} con un TTL de 10 minutos. La página del Optimizador, al cargarse, detecta esa recomendación y la muestra como banner informativo, además de pre-seleccionar las plantillas adecuadas al perfil.

% =============================================================================
\section{X-Ray: análisis de la cartera}
% =============================================================================

La página \texttt{X-Ray} desagrega cualquier cartera guardada según múltiples dimensiones, al estilo del informe que ofrece Morningstar para sus fondos.

\subsection{Distribución por clase de activo}

Detección automática mediante una mezcla de:
\begin{enumerate}
  \item Mapeo explícito de tickers populares (\textit{BND}, \textit{TLT}, \textit{GLD}, etc. → su clase correspondiente). Es el método más fiable.
  \item \texttt{quoteType} de yfinance (CRYPTOCURRENCY → Alternatives, etc.).
  \item Búsqueda de palabras clave en nombre, sector, industria y categoría (\textit{bond}, \textit{treasury}, \textit{REIT}, \textit{gold}, etc.).
  \item Por defecto: Equity para EQUITY/ETF/MUTUALFUND.
\end{enumerate}

Se muestran seis clases: Equity, Bonds, Real Estate, Commodities, Cash, Alternatives.

\subsection{Sectores y geografía}

Sectores: agregación de pesos por sector de yfinance (Technology, Healthcare, etc.) en un \textit{donut chart}.

Geografía: mapa coroplético del mundo (\texttt{react-simple-maps} con topojson) coloreado por intensidad de exposición. Se mantiene un mapeo de nombres de países a códigos ISO numéricos para el matching.

\subsection{Style Box Morningstar}

Matriz $3\times 3$ que cruza capitalización (Large/Mid/Small) y estilo de inversión (Value/Blend/Growth):

\begin{description}[leftmargin=1.4cm]
  \item[Capitalización] $> \$10$B: Large; $\$2$B--$\$10$B: Mid; $< \$2$B: Small.
  \item[Estilo] Heurística sobre P/E y P/B: scoring con $\pm 1$ por tramo, $\le -1$ Value, $\ge +1$ Growth, intermedio Blend.
\end{description}

Cada celda se colorea con intensidad proporcional al peso de la cartera que cae en esa combinación.

\subsection{Tipo de activo Morningstar}

Clasificación por sector en tres grandes grupos: Cíclico (Materials, Financials, Consumer Cyclical, Real Estate), Sensible (Communication Services, Energy, Industrials, Technology) y Defensivo (Consumer Defensive, Healthcare, Utilities).

\subsection{Análisis de dividendos}

\begin{itemize}
  \item \textbf{Yield ponderado de la cartera}: $\sum_i w_i \cdot y_i$. Se traduce a "ingresos anuales aproximados sobre 10\,000\,€".
  \item \textbf{Frecuencia de pago}: detectada del histórico real de pagos (no del campo declarado por yfinance, que a veces es incorrecto). Diferenciamos mensual, trimestral, semestral y anual según la mediana de días entre pagos.
  \item \textbf{Top 3 contribuyentes} al yield total con barras de aporte en puntos porcentuales.
\end{itemize}

\subsection{Historial de versiones (snapshots)}

Cada vez que el usuario modifica los pesos de una cartera guardada, el estado anterior se persiste como \textit{snapshot}. El X-Ray ofrece un modal con un \textit{timeline} vertical donde se muestra cada versión, con un \textit{diff} respecto a la siguiente que indica activos añadidos, eliminados y modificados (con la variación en puntos porcentuales).

% =============================================================================
\section{Comparador de carteras}
% =============================================================================

La página \texttt{Comparador} permite seleccionar hasta 4 carteras guardadas y ejecutar un backtest paralelo (\texttt{Promise.all}) sobre el mismo periodo. El resultado es:

\begin{itemize}
  \item Curvas de \textit{equity} superpuestas en una sola gráfica (todas con base 100), con benchmark SPY como referencia.
  \item Tabla comparada con 8 métricas (Rentabilidad acumulada, Retorno anualizado, Volatilidad, Sharpe, Sortino, Calmar, Max DD, Beta).
  \item Composición lado a lado con barras horizontales, que permite ver concentraciones y solapamientos de un vistazo.
\end{itemize}

% =============================================================================
\section{Plantillas de cartera}
% =============================================================================

En \texttt{frontend/src/portfolioTemplates.ts} hay 13 plantillas preconfiguradas inspiradas en estrategias famosas:

\begin{description}[leftmargin=2cm]
  \item[Conservadoras] Defensiva con Dividendos · 60/40 Global · Ingresos \& Utilities · Bond Ladder (solo renta fija) · Permanent Portfolio (Harry Browne).
  \item[Moderadas] Core S\&P 500 + Internacional · Sectores Balanceados · Blue Chips US · All Weather (Ray Dalio) · Income con bonos + dividendos.
  \item[Agresivas] Mega-cap Tech · Disrupción e Innovación · Semis + Emergentes.
\end{description}

El optimizador filtra las plantillas según el perfil detectado del usuario y permite ver todas con un toggle.

% =============================================================================
\section{Capa API y autenticación}
% =============================================================================

\subsection{Endpoints principales}

\begin{description}[leftmargin=1.6cm]
  \item[\texttt{POST /auth/register}] Registro con email + password (validación, hash bcrypt).
  \item[\texttt{POST /auth/login}] Login \texttt{OAuth2PasswordBearer}, devuelve JWT.
  \item[\texttt{GET /auth/profile}] Perfil del usuario autenticado.
  \item[\texttt{PUT /auth/profile}] Actualización de capital base y tolerancia al riesgo.
  \item[\texttt{GET /assets/search?q=...}] Búsqueda de tickers en Yahoo Finance (cacheada).
  \item[\texttt{GET /assets/\{ticker\}/info}] Metadatos del activo: sector, país, capitalización, dividend yield, asset class.
  \item[\texttt{POST /portfolio/optimize}] Optimización con uno de los 5 métodos.
  \item[\texttt{POST /portfolio/monte-carlo}] Simulación Monte Carlo de la cartera.
  \item[\texttt{POST /portfolio/fiscalidad}] Análisis fiscal IRPF.
  \item[\texttt{POST /backtesting/run}] Backtest histórico con rebalanceo y comisiones.
  \item[\texttt{POST /portfolio/}] Guardar cartera del usuario autenticado.
  \item[\texttt{GET /portfolio/}] Listar carteras del usuario.
  \item[\texttt{PUT /portfolio/\{id\}}] Actualizar cartera (crea snapshot del estado anterior).
  \item[\texttt{DELETE /portfolio/\{id\}}] Eliminar cartera (cascada sobre activos y snapshots).
  \item[\texttt{GET /portfolio/\{id\}/snapshots}] Historial de versiones.
\end{description}

\subsection{Modelos de base de datos}

Cinco tablas, definidas en \texttt{backend/app/db/models.py}:

\begin{itemize}
  \item \texttt{Usuario}: id, email, password\_hash, capital\_base, tolerancia\_riesgo.
  \item \texttt{Cartera}: id, usuario\_id, nombre\_estrategia, fecha\_creacion.
  \item \texttt{Activo}: id, isin\_ticker, nombre\_fondo. Catálogo de activos referenciados.
  \item \texttt{ActivoCartera}: tabla N:N entre Cartera y Activo, con peso\_asignado.
  \item \texttt{CarteraSnapshot}: id, cartera\_id, fecha, nombre, pesos (JSON), motivo.
\end{itemize}

Las migraciones están versionadas con Alembic en \texttt{backend/alembic/versions/}.

\subsection{Seguridad}

\begin{itemize}
  \item Contraseñas hasheadas con bcrypt (factor de coste configurable).
  \item JWT firmado con \texttt{SECRET\_KEY} y expiración configurable (8 horas por defecto, demasiado largo para producción real).
  \item Interceptor del frontend que cierra sesión automáticamente ante un 401 fuera de los endpoints de auth.
  \item Validación de pesos: ahora rechaza valores negativos (no se permiten posiciones cortas) y comprueba suma = 1.
\end{itemize}

% =============================================================================
\section{Decisiones de diseño relevantes}
% =============================================================================

\begin{enumerate}
  \item \textbf{Bootstrap histórico vs paramétrico Gaussiano} para Monte Carlo (decisión razonada en la sección correspondiente).
  \item \textbf{Multi-start SLSQP} con semilla fija para reproducibilidad. Alternativas más sofisticadas (CVXPY, IPOPT) requerirían dependencias pesadas y aportarían poco para los tamaños de problema típicos.
  \item \textbf{Caché en memoria con TTL} en lugar de Redis. Suficiente para la escala del TFG y elimina una dependencia externa.
  \item \textbf{\texttt{auto\_adjust=True}} en yfinance: los precios devueltos son total return. Decisión documentada explícitamente al usuario.
  \item \textbf{No tabla de transacciones individuales}: el coste fiscal se modela paramétricamente. Para un asesor real haría falta historial de operaciones; para un TFG la simplificación es aceptable y declarada.
  \item \textbf{Lógica del planner en TypeScript puro} (no en backend): permite recalcular en tiempo real sin viajes al servidor y mantiene la lógica testeable de forma independiente.
\end{enumerate}

% =============================================================================
\section{Limitaciones reconocidas}
% =============================================================================

\begin{enumerate}
  \item Los retornos esperados $\mu$ se estiman como media histórica simple, lo cual es notoriamente ruidoso. Mejoras posibles: \textit{shrinkage} (Ledoit-Wolf), Black-Litterman.
  \item La matriz de covarianzas tampoco se regulariza. Funciona bien para 5--10 activos pero degrada con $N$ grande.
  \item No se modela slippage ni \textit{market impact}.
  \item El bootstrap no genera escenarios fuera de la muestra histórica.
  \item La fiscalidad simplificada no contempla compensación de pérdidas, casos transfronterizos ni reglas anti-recompra.
  \item No hay deploy en producción todavía: la app corre solo en local con Docker Compose.
\end{enumerate}

% =============================================================================
\section{Referencias bibliográficas}
% =============================================================================

\begin{thebibliography}{99}

\bibitem{markowitz1952}
Markowitz, H. (1952). \textit{Portfolio Selection}. Journal of Finance, 7(1), 77--91.

\bibitem{demiguel2009}
DeMiguel, V., Garlappi, L., Uppal, R. (2009). \textit{Optimal Versus Naive Diversification: How Inefficient is the 1/N Portfolio Strategy?}. Review of Financial Studies, 22(5), 1915--1953.

\bibitem{maillard2010}
Maillard, S., Roncalli, T., Teïletche, J. (2010). \textit{The Properties of Equally Weighted Risk Contribution Portfolios}. Journal of Portfolio Management, 36(4), 60--70.

\bibitem{choueifaty2008}
Choueifaty, Y., Coignard, Y. (2008). \textit{Toward Maximum Diversification}. Journal of Portfolio Management, 35(1), 40--51.

\bibitem{rockafellar2000}
Rockafellar, R. T., Uryasev, S. (2000). \textit{Optimization of Conditional Value-at-Risk}. Journal of Risk, 2(3), 21--41.

\bibitem{artzner1999}
Artzner, P., Delbaen, F., Eber, J. M., Heath, D. (1999). \textit{Coherent Measures of Risk}. Mathematical Finance, 9(3), 203--228.

\bibitem{browne1999}
Browne, H. (1999). \textit{Fail-Safe Investing: Lifelong Financial Security in 30 Minutes}. St. Martin's Griffin.

\bibitem{sortino1994}
Sortino, F., Price, L. (1994). \textit{Performance Measurement in a Downside Risk Framework}. Journal of Investing, 3(3), 59--64.

\bibitem{ley35}
\textit{Ley 35/2006, de 28 de noviembre, del Impuesto sobre la Renta de las Personas Físicas} (BOE-A-2006-20764, con sus modificaciones posteriores).

\end{thebibliography}

\end{document}
